\subsection{Il modello di output e la classificazione dell'accuratezza}
\label{sec:output-model}

Per consentire allo scheduler di confrontare le varianti in termini di 
accuratezza, ciascuna variante è dotata di un \texttt{OutputModel} che 
descrive la qualità del suo output. Il modello supporta due tipologie 
di classificazione, scelte in base alla natura computazionale della 
funzione:

\begin{itemize}
    \item \textbf{Tipo \texttt{error}}: adatto a funzioni il cui output 
    è un valore numerico misurabile. L'accuratezza è espressa tramite 
    il campo \texttt{error\_estimate}, che rappresenta l'errore relativo 
    atteso rispetto alla variante di riferimento. Una variante esatta 
    presenta \texttt{error\_estimate = 0.0}, mentre varianti approssimate 
    presentano valori positivi proporzionali alla deviazione introdotta 
    dalla semplificazione algoritmica adottata.

    \item \textbf{Tipo \texttt{quality}}: adatto a funzioni il cui 
    output non è facilmente quantificabile con un errore numerico, 
    come classificatori o modelli di machine learning. L'accuratezza 
    è espressa tramite un livello discreto: \texttt{high}, 
    \texttt{medium} o \texttt{low}, che riflette la qualità complessiva 
    del risultato prodotto dalla variante.
\end{itemize}

Il ~\autoref{lst:imageclassification-json} mostra un esempio 
applicativo del tipo \texttt{quality} per la funzione 
\texttt{ImageClassification}, in cui la variante base utilizza un 
modello di rete neurale più accurato ma più oneroso, mentre la 
variante \texttt{Light} adotta un modello più compatto a qualità ridotta.

\begin{lstlisting}[
    language=json,
    caption={File di definizione delle varianti per la funzione 
    \texttt{ImageClassification}. Il campo \texttt{quality} esprime 
    l'accuratezza in forma discreta, adatta a output non numericamente 
    quantificabili.},
    label={lst:imageclassification-json}
]
{
  "variants": [
    {
      "id": "base",
      "runtime": "python-ml",
      "entry_point": "ImageClassification.handler",
      "src": "ImageClassification.py",
      "energy": {
        "cold_start_joule": 1.894736842,
        "invocation_joule": 1.455263158
      },
      "output": { "type": "quality", "quality": "high" },
      "is_approximate": false
    },
    {
      "id": "Light",
      "runtime": "python-ml",
      "entry_point": "ImageClassificationLight.handler",
      "src": "ImageClassificationLight.py",
      "energy": {
        "cold_start_joule": 1.071052329,
        "invocation_joule": 0.22894768
      },
      "output": { "type": "quality", "quality": "low" },
      "is_approximate": true
    }
  ]
}
\end{lstlisting}

Lo scheduler converte internamente entrambe le rappresentazioni in un 
punteggio numerico (\emph{ErrorScore}) secondo una scala uniforme: 
per il tipo \texttt{error} viene utilizzato direttamente il valore di 
\texttt{error\_estimate}, mentre per il tipo \texttt{quality} i livelli 
\texttt{high}, \texttt{medium} e \texttt{low} vengono mappati 
rispettivamente a $0.0$, $1.0$ e $2.0$. La selezione privilegia 
sempre la variante con il punteggio più basso compatibile con il 
vincolo energetico, garantendo che la qualità del risultato sia 
massimizzata nei limiti imposti dalla richiesta.

La Tabella~\ref{tab:output-models} riassume il modello di output 
adottato per ciascuna delle funzioni utilizzate negli esperimenti.

\begin{table}[h]
    \centering
    \begin{tabular}{llll}
        \toprule
        \textbf{Funzione} & \textbf{Tipo} & 
        \textbf{Variante base} & \textbf{Variante approssimata} \\
        \midrule
        Fibonacci         & \texttt{error}   & 0.0  & 0.0 \\
        PiLeibniz         & \texttt{error}   & 0.0  & $1.99 \times 10^{-8}$ \\
        Sqrt              & \texttt{error}   & 0.0  & 0.043 \\
        VectorMagnitude   & \texttt{error}   & 0.0  & 0.04 \\
        SentimentAnalysis & \texttt{quality} & high & low \\
        ImageClassification & \texttt{quality} & high & low \\
        \bottomrule
    \end{tabular}
    \caption{Modello di output adottato per ciascuna funzione. Per il 
    tipo \texttt{error} il valore riportato è \texttt{error\_estimate}; 
    per il tipo \texttt{quality} il livello discreto assegnato.}
    \label{tab:output-models}
\end{table}